../cictikz/src/cictikz/data/tex/cictikz_preamble.tex

% The grounded gate NMOS on its own, with the zap current forced into it.
%
% Gate, source and bulk on ground, and 2.6 A into the drain. The whole
% point of the figure is that pair of facts side by side: every terminal
% you can name is at zero, so the device is off by any model you have
% been taught, and yet it is expected to sink an amp scale current. What
% carries it is the parasitic lateral bipolar, which the BSIM model in
% front of you does not have - that is the sentence the chapter spends
% the next two pages on.
%
% 2.6 A is the HBM number from earlier in the chapter: 4 kV across the
% 1.5 k body resistance.

\begin{circuitikz}[american, thick, transform shape, circuit ee IEC]

  ../cictikz/src/cictikz/data/tex/cictikz_lib.tex

  %- transistor, source on ground
  \draw (0,0) \vground;
  \draw (0,0) \lvnmos{M1}{};

  %- gate to its own ground, beside the source: drawn as a separate
  %- symbol rather than a wire back to the source pin, because the
  %- chapter's next figure is a cross section where the gate contact and
  %- the source contact really are two different pieces of metal
  \draw (M1.gate) -- (-1.1,0.8);
  \draw (-1.1,0.8) -- (-1.1,0);
  \draw (-1.1,0) \vground;

  %- the zap, into the drain
  \draw (0,{2*\grid}) to [I] (0,\grid);
  \draw (0,{2*\grid}) -- (0,{2*\grid+0.6});
  \node[red, anchor=west] at (0.55,{1.5*\grid}) {$2.6$ A};

\end{circuitikz}
\end{document}
